\section{有效相互作用与穿衣顶点}

Gross 第 28 章：在相互作用线上插入全部极化图得穿衣相互作用 $W$；再引入不可约顶点 $\Lambda$，并强调避免双重计数。

\subsection{极化插入与 $W$}

对裸线 $v$ 反复插入极化泡，几何级数给出
\begin{equation}
  W(1,2)
  =v(1,2)
  +\int\mathrm{d}3\,\mathrm{d}4\,
  v(1,3)\,
  \Pi(3,4)\,
  W(4,2),
\label{eq:W_def}
\end{equation}
或在均匀体系
\begin{equation}
  W(\mathbf{q},\omega)
  =\frac{v_{\mathbf{q}}}{1-v_{\mathbf{q}}\Pi(\mathbf{q},\omega)}
  =\frac{v_{\mathbf{q}}}{\varepsilon(\mathbf{q},\omega)}.
\end{equation}
RPA 取 $\Pi=\Pi_0$（Lindhard）。等离激元是 $\varepsilon(\mathbf{q},\omega)=0$ 的极点，因而出现在 $W$ 与密度响应上，\textbf{不是}单粒子 $G$ 的极点。

\subsection{不可约顶点 $\Lambda$}

\textbf{顶点插入}：接在相互作用端点与两条 $G$ 腿之间、且不能再切开自能或极化泡的图。不可约顶点 $\Lambda(k;q)$ 是全部此类插入之和。示意
\begin{equation}
  \Lambda
  =1
  +I\,G G\,\Lambda,
\end{equation}
$I=\delta\Sigma/\delta G$ 时与守恒近似一致。

\subsection{双重计数警告}

把 $\Lambda$ 写入自能时，\textbf{只能挂在有效相互作用的一端}；若两端都挂 $\Lambda$，会与已含在 $W$ 或 $\Sigma$ 中的图重复（Gross 图 28.12--28.13）。正确的示意结构是
\begin{equation}
  \Sigma
  \sim
  G\,W\,\Lambda
  \quad\text{（一端顶点）},
\end{equation}
而非 $\Lambda G W\Lambda$。

\subsection{部分求和的选择}

\begin{itemize}
  \item \textbf{环图}：长程库仑、等离激元、GMB；
  \item \textbf{梯子图}：短程强排斥、低密气、顶点修正主导；
  \item \textbf{骨架自洽}：$\Sigma[G]$、$W[\Pi]$、$\Lambda$ 联立迭代。
\end{itemize}
何种最优取决于力程与能标。Ward 恒等式把 $q\to0$ 的 $\Lambda$ 与 $\partial G^{-1}/\partial\mu$ 相连，保证电荷守恒，并衔接微观 Landau 参数。
